\scanpage{13}
% 1–7
et $(P^\varepsilon)'(z)=Nz^{N-1}+(N-1)a_1z^{N-2}\ldots+a_{N-1}$ soit $\in k^*$ (donc en fait $\in k_0^*$), de sorte qu'il existe un unique $F\in S_n$ telle que $P(F)=0$ et $\varepsilon(F)=z$. Sous ces conditions, on a $F\in A_n$. [En d'autres termes, le couple $(A_n,\mu_n=\mathfrak m_n\cap A_n)$ est (non seulement zariskien, selon le cor.~1, mais même) hensélien]

% 8–11
\deleted{En effet, \inserted{\unsure{quitte à remplacer|quitte à rempl.} $Z$ par $Z\unsure{+|-}z$, sps $z=0$, donc $\unsure{a_N|f_N}=\varepsilon(f_N)=0$} on regarde $X_{n+1}\to X_{n+1}$ donné par}
\[
\deleted{(x_1,\ldots,x_n,x_{n+1})\longmapsto (x_1,\ldots,x_n,x_{n+1}^{N+1}+f_1(x_1,\ldots,x_n)x_{n+1}^{N}\ldots)}
\]
\deleted{L'hyp. sur}

Plus généralement

% 12–16
\uline{Cor.~3} Soit $P\in A_{n+1}[t_1,\ldots,t_n,Z]$, $P=f_0+f_1Z+f_2Z^2+\ldots$ [les $f_i\in S_n(t_1,\ldots,t_n)$] supposons $\varepsilon(f_0)=0$ i.e. $\varepsilon(P)=0$, et $\varepsilon(f_1)\in k^*$ i.e. $\varepsilon(\frac{\partial P}{\partial Z})\in k^*$, de sorte qu'il existe une unique $F\in S_n$ telle que $P(t_1,\ldots,t_n,F(t_1,\ldots,t_n))=0$ et $\varepsilon(F)=0$. Alors $F\in A_n$.

% 17–24
Soit en effet $u:X_{n+1}\to X_{n+1}$ donné par $(x_1,\ldots,x_n,z)\longmapsto(x_1,\ldots,x_n,P(x_1,\ldots,x_n,z))$, alors l'hyp. dit que c'est compatible avec les ponctuations, et que c'est un isomorphisme formel. Comme l'hyp. $P\in A_{n+1}$ implique $u$ un $A_*$-morphisme, $u$ est un $A_*$ isomorphisme et par suite $u^{-1}$ est un $A_*$-morphisme. Donc
\[
\unsure{u^{-1*}|u^{-1}}(x_1,\ldots,x_n,z)=(x_1,\ldots,x_n,\deleted{\unsure{z}}\varphi(x_1,\ldots,x_n,z))
\]
avec $\varphi\in A_{n+1}(t_1,\ldots,t_n,t_{n+1})$, donc $\inserted{F=}\varphi(t_1,\ldots,t_n,0)\in A_n(t_1,\ldots,t_n)$

% 17–24, left
% margin note written sideways
\margin{a) b) c) d) impliquent-ils que $S_n$ soit fid. plat sur $A_n$ (si $k$ \deleted{\unsure{l'|le}}l'est \deleted{\unsure{plat|fini}} sur $k_0$) ?}

% 25–27
\deleted{d) (\unsure{Condition différentielle|Conditions différentielles}) Les $A_n$ sont \uline{stables} par les opérateurs différentiels standard}
\[
\deleted{D_{\unsure{p_i|p|p_j}}=D_{(p_1,\ldots,p_n)}\quad\left(=\frac{1}{p_1!\ldots p_n!}\frac{\partial^{|p|}}{\partial x_1^{p_1}\ldots\partial x_n^{p_n}}\right)}
\]
% archive sheet number: 12
